\documentclass[11pt]{article}
\usepackage[margin=.8in]{geometry}
\usepackage{amsmath,amssymb}
\usepackage[colorlinks=true,urlcolor=blue]{hyperref}
\hypersetup{pdftitle={Reading guide: one-edge fractional graph response}}
\setlength{\parindent}{0pt}
\setlength{\parskip}{5pt}
\begin{document}
\begin{center}
{\Large One-edge fractional graph response}\par\vspace{3pt}
{\large Results and reading guide}\par\vspace{3pt}
November 15, 2026
\end{center}

\textbf{The problem.}
For a connected weighted graph with unit vertex masses, increase or
insert one edge: $L'=L+d bb^T$, $d>0$, $b=e_p-e_q$.
For fixed $0<\alpha<1$, measure the response in the original energy:
\[
 C=(L^\dagger)^{\alpha/2}(L'^\alpha-L^\alpha)(L^\dagger)^{\alpha/2},
 \qquad \tau=d b^TL^\dagger b.
\]
Inverse powers vanish on constants. Let $N_\varepsilon(C)$ count
eigenvalues strictly greater than $\varepsilon$. This is the optimal
rank for relative spectral truncation to tolerance $\varepsilon$.

\textbf{Exact unrestricted count.}
Positivity and the determinant lemma give
\[
 C\succeq0,\quad
 \sum_j\log(1+\lambda_j(C))=\alpha\log(1+\tau),\quad
 N_\varepsilon(C)\le
 M_\tau(\varepsilon):=
 \max\left\{0,\left\lceil
 \frac{\alpha\log(1+\tau)}{\log(1+\varepsilon)}
 \right\rceil-1\right\}.
\]
The strict ceiling-minus-one handles integral thresholds. The upper
bound is the exact worst case on weighted paths with ordinary unit
vertex masses (Theorem 3.1). Its leading term is
$\alpha\log(1+\tau)/\varepsilon$.

\textbf{Unweighted universality.}
For $\tau=1$, the same exact count already occurs between simple,
unweighted, connected, bipartite graphs differing by insertion of one
absent edge. The inserted endpoints have resistance \emph{exactly one}
in every finite host (Theorem 4.1). Every finite list $q_i\ge1$ with
$\prod_iq_i=2^\alpha$ is a limit of relative spectra, with additional
correction eigenvalues tending to zero. Graph size and degree may grow.

\textbf{Three different claims.}
In fixed matrix dimension, spectral closure requires equality of the
determinant budget. For a finite correction target $q_i-1$ padded with zeros while graph
dimension grows, uniform closure permits $\prod_iq_i\le(1+\tau)^\alpha$:
many tiny entries can carry the deficit. Full-budget constructions can
instead converge in $\ell^1$. Exact finite matrix realization imposes
pairwise distinct values above one; repeated positive corrections are
limit statements (Appendix B).

\textbf{Degree changes the sharp rank law.}
Among simple unweighted hosts, for doubling one existing unit edge
(Section 5), paths and cycles have
worst-case count
$W_{\mathrm{path},\alpha}(\varepsilon)\sim W_{2,\alpha}(\varepsilon)
\sim\log^2(1/\varepsilon)/(2\pi^2)$. For every fixed integer $d\ge3$,
Section 5 proves
\[
 \lim_{\varepsilon\downarrow0}\varepsilon W_{d,\alpha}(\varepsilon)
 =A_{d,\alpha}:=\alpha\log\frac{2d}{d+2}.
\]
The same sharp coefficient is attained by bipartite hosts and bridge
edits. In particular $A_{3,\alpha}=\alpha\log(6/5)$.
The upper proof subtracts the entire rooted local response; its wired
resistance and the finite graph's free-resistance bound are different
quantities. Multiport reservoirs approach the remaining budget with
many nearly equal modes. These are asymptotic coefficients, not exact
finite-threshold counts. The degree-two upper compares every finite grid
with one fixed continuum Cauchy operator; it is uniform over all sizes
and edge locations. The centered-path limit supplies the matching lower. Doubling changes the edited edge's conductance from one to two.

\newpage
\textbf{A charged construction from the graph.}
Theorem 6.1 returns an explicit $X\in\mathbb R^{n\times r}$, with
probability at least $1-\delta$, satisfying
\[
 -\frac\varepsilon8L^\alpha
 \preceq L'^\alpha-L^\alpha-XX^T
 \preceq\frac{3\varepsilon}4L^\alpha,
 \qquad
 r<\frac{\alpha\log(1+\tau)}{\log(1+\varepsilon/4)}\quad(r>0).
\]
For simple unweighted hosts with one existing edge doubled,
Corollary 6.6 attains the sharp leading rank at every fixed $d\ge2$:
\[
 r\le
 \begin{cases}
 (1/(2\pi^2)+o_\alpha(1))\log^2(1/\varepsilon),&d=2,\\
 (A_{d,\alpha}+o_{d,\alpha}(1))/\varepsilon,&d\ge3.
 \end{cases}
\]
For $d=2$, a fixed cutoff parameter $\rho=1/2$ suffices. For $d\ge3$,
use $\rho=1/k$, where $k$ is the least integer with $2^k\ge4/\varepsilon$,
found by repeated doubling.
Here the error guarantee is relative operator norm at most $\varepsilon$;
rank and error hold on the same success event. The remainder is uniform
over the graph class, without an effective rate. The sharp bridge lower
examples force this constant for any rank approximant, even without
positivity.

The expected arithmetic work is $\widetilde O_\alpha(m+n)$.
Suppressed factors are polynomial in logarithms of
$n,1+\tau,1/\varepsilon,1/\delta$; they contain no graph gap,
spectral norm, or conductance ratio. Both raw and Gram quadrature widths
are $O_\alpha(\log^2(2+n(1+\tau)/\varepsilon))$.
The proof charges feasible electrical flows, all shifted SDD solves,
small-Gram formation, Jacobi rotations, and every factor entry.
Certified primal/dual resistance intervals and explicit work guards
control expected work even on failed solver calls.

\textbf{Model and scope.}
This is a positive-edit, unit-vertex-mass, exact-real arithmetic theorem.
General fixed-$\alpha$ scalar powers are scalar operations; the
square-root case has an arithmetic-and-square-root implementation.
It is not a finite-word stability result. The computed factor has
\emph{two-sided} relative error; it need not lie below the true update.
The algorithm may reconstruct solver data per edit. It does not prove a
sublinear maintained dynamic data structure.

If $P=L^\alpha+XX^T$, then
$(1-\varepsilon)P\preceq L'^\alpha\preceq(1+\varepsilon)P$.
Applying $P^{-1}$ still requires a separately charged baseline
$L^{-\alpha}$ action. The theorem constructs the correction and proves
its quality; it does not make that baseline action free.

\textbf{Where to read.}
The complete manuscript is
\href{single-edge-response.pdf}{\texttt{single-edge-response.pdf}}.
Its proof routes are:

\begin{tabular}{@{}p{.76\textwidth}p{.18\textwidth}@{}}
Question & Sections\\ \hline
Budget, metric, exact rank law & 1--2\\
Weighted-path realization & 3\\
Actual unweighted insertion & 4\\
Degree laws, sharp coefficient, local-limit rigidity & 5\\
Explicit construction and all arithmetic costs & 6\\
Deletion, topology, comparison with prior methods & 7\\
General Araki--Horn matrix range and batch envelope & Appendix A\\
Exact finite rank-one attainment and simplicity & Appendix B
\end{tabular}

\textbf{Classical inputs and status.}
Araki, Weyl--Horn/Kostant, Schur asymptotics, total positivity, resistor
laws, coarse graining, Stahl approximation, Laptev/Widom asymptotics,
expander existence, Fredholm determinants, free/wired networks and Kron
reduction, rational matrix-function updates, sinc/contour
quadrature, and nearly linear electrical-flow solvers are credited in
the manuscript. Internal independent reviews found no remaining
mathematical blocker in the integrated claims. These reviews are not
formal or external peer review, and do not certify historical novelty.
Efficient extremal graph sizes, effective higher-degree asymptotic
rates, and finite-word stability remain open here.
\end{document}
